\section{Conflicting modalities require correctness evidence and a separate causal test}
\label{exec23:conflict-scope}

A bank analyst may encounter a chart, its caption and a service narrative that
support different answers. Zhu and colleagues study an adjacent problem in
\emph{Unraveling Cross-Modality Knowledge Conflicts in Large Vision-Language
Models}. The inspected original is arXiv version 1, dated 4 October 2024;
the archive's later version 2 is a separate edition. Their text and image
question-answering experiments motivate checks for inconsistent representations,
but do not establish causal identification or bank decision performance.
\citep{exec23_cross_modality_conflict}

Let $T$ and $V$ be textual and visual answer labels for the same task and let
$G$ be its gold label. Define flip rate $F=P(T\ne V)$ and accuracies
$a_T=P(T=G)$ and $a_V=P(V=G)$. The paper proposes a lower bound on its
knowledge-conflict rate of the form $\mathrm{CR}\geq F-(a_T-a_V)$, attributing
a net accuracy gap to performance-gap samples. That attribution is an
additional premise about unobserved knowledge, rather than an identified
consequence of observing predictions.~\citep{exec23_cross_modality_conflict}

Partition the observations by the correctness of each answer. Cases where
text alone is correct and cases where vision alone is correct both flip.
Both-correct cases do not flip; both-wrong cases may flip. Subtracting the
accuracy difference gives the exact prediction identity
\begin{equation}
 F-(a_T-a_V)=
 2P(T\ne G,V=G)+P(T\ne G,V\ne G,T\ne V).
 \label{exec23:conflict-identity}
\end{equation}
This is a project derivation about answer labels. It supplies no observation
of stored knowledge and no causal explanation of the errors. In a four-item
authored example with gold answers AAAA, text AABB and vision BBAA, each
accuracy is one half, yet every answer flips and the right side equals one.
With text and vision both BBBB, agreement is perfect and both accuracies are
zero. Agreement, correctness and causal validity therefore need distinct
evaluation records.~\citep{exec23_context_reference}

\subsection{The decoding proposal and what the printed results establish}
\label{exec23:conflict-decoding}

The paper reports that confidence, confidence shifts and minimum-variance
strategies did not reliably choose correct answers in the tested conflicts.
Its dynamic contrastive decoding (DCD) uses answer confidence to scale
log-probabilities and subtracts the less confident modality from the more
confident modality. It requires access to the relevant probabilities and
both modality paths. It is a proposed inference-time answer rule, not a
causal intervention on the bank's data-generating process. In particular,
the earlier finding about confidence limits qualifies the heuristic used
to select the subtraction direction.~\citep{exec23_cross_modality_conflict}

Table 4 reports LLaVA-34b accuracy rising from 80.99 to 83.35 on ViQuAE and
from 66.02 to 68.14 on InfoSeek: gains of 2.36 and 2.12 percentage points,
averaging 2.24 points. These are author-reported cells, not reproduced model
runs. The same table contains an InstructBlip-7b InfoSeek tie at 50.53, so
the printed claim of strict improvement across every case is too broad.
There is also a baseline inconsistency: Table 2 prints 82.46 for LLaVA-34b
ViQuAE textual accuracy, whereas Table 4 prints 80.99. Neither value should
silently replace the other.~\citep{exec23_cross_modality_conflict,exec23_context_reference}

Selected Table 5 cells require a stronger arithmetic qualification. Against
the displayed visual baseline accuracy 53.26, the reminder and answer-conflict
values 53.99 and 54.58 differ by $+0.73$ and $+1.32$ percentage points;
the table instead prints $-1.66$ and $-1.07$. Against the displayed recognized
accuracy 58.11, values 57.25 and 58.51 differ by $-0.86$ and $+0.40$;
the printed changes are $-2.53$ and $-1.27$. These four inconsistencies do
not resolve which underlying experiment or baseline is intended. The selected
InfoSeek accuracy changes from 22.11 to 21.25 and 20.23 do match the printed
$-0.86$ and $-1.88$. Preserve the raw cells and their edition while seeking
data or clarification; do not use the inconsistent signs to rank a mitigation.
\citep{exec23_cross_modality_conflict,exec23_context_reference}

\subsection{A bounded route from this literature to a banking prototype}
\label{exec23:conflict-prototype}

The paper's public repository is a candidate for inspection. Its code, data,
license, environment and raw results have not been acquired or qualified here.
The proposal can nevertheless specify useful tests using public or authored
materials. Construct matched chart, table and narrative versions of the same
known-law problem, deliberately introduce a unit change, stale date, swapped
axis or unsupported causal caption, and retain a source-backed gold answer
for each representation. Evaluate answer correctness and conflict detection
separately; count abstentions and incorrect agreement. Require a causal
question to pass the separate target, graph and identification checks developed
above even when all representations agree.

Compare a single-modality baseline, source inspection, explicit disagreement
reporting and, where logits and rights permit, a reproduced DCD rule. Charge
for both modality passes, repeated confidence samples and human adjudication.
The inspected appendix's ten prompt or dropout samples concern its confidence
analysis; they are not a general guarantee about uncertainty or an automatic
ten-fold price for DCD. Freeze tasks, corruptions, costs and denominators before
evaluation. A generated monetary target can expose a wrong decision under a
declared law, while authentic bank utility remains a later experiment. This
sequence makes the expensive study conditional on a demonstrated prototype merit.
